\documentclass[12pt]{article}

\usepackage[T1]{fontenc}
\usepackage[utf8]{inputenc}
\usepackage[margin=1in]{geometry}
\usepackage{setspace}
\usepackage{parskip}
\usepackage{titlesec}
\usepackage[hidelinks]{hyperref}

\titleformat{\section}{\normalfont\bfseries\large}{}{0pt}{}
\titlespacing*{\section}{0pt}{1.4em}{0.6em}

\setstretch{1.15}

\begin{document}

\begin{center}
	{\Large\bfseries Primary Source Worksheet \#4}\\[0.6em]
	RELI 252 \quad\textbullet\quad Dr.\ Clark\\[0.3em]
	Greyson Knippel
\end{center}

\vspace{1em}
\hrule
\vspace{1.5em}

\noindent\textbf{Primary Source Title:} \textit{The Travail of a Female Colored Evangelist}

\vspace{0.5em}

\noindent\textbf{Primary Source Author:} Amanda Berry Smith

\section*{What is the author's intent with this piece?}

Smith wants her readers to see that she was shut out everywhere and kept preaching anyway. She is trying to show that her authority is real even though no church ever gave it to her.

She does this by telling on people without ever attacking them. When a woman throws her out of a church in New York, Smith records the exact words, ``Pass out, pass out,'' and then says she prayed for the woman. When two ladies in Philadelphia tell her the meeting is for wealthy guests and ``thee had better not go to-day,'' Smith just writes ``I never said a word.'' She lets the scene do the damage. This keeps her looking humble while the reader does the judging.

She also wants to prove she never forced herself anywhere. She says it directly: ``I don't think that anybody can ever say that Amanda Smith pushed herself in where she was not wanted.'' That matters because the main charge against women preachers was that they were out of place. She is answering it before anyone makes it.

The biggest thing she wants the reader to accept is where her authority comes from. When AME ministers assume she came to Nashville to demand ordination, she says the thought never crossed her mind, because ``I had received my ordination from Him.'' She is claiming that a call from God outranks a vote by a church body.

She also wants the reader to notice things changed. She points out that people who once thought she would ruin a meeting now think the meeting fails without her.

\section*{Compare and contrast with a previously assigned primary source}

Smith's autobiography and Jarena Lee's make almost the same argument fifty years apart, which is pretty insteresting when reading them together.

Both women were Black, both were Methodist, both felt called to preach, and both were blocked by the AME. Both ended up claiming authority directly from God because no institution would give it to them. Lee asks whether Christ is ``a whole Saviour, instead of a half one'' if women cannot preach. Smith says she already has her ordination from the one who said ``I have chosen you, and ordained you.'' Same situation. Both of them did not simply wait for permission .

The difference is what they were up against. Lee ran into a rule. Richard Allen told her the Discipline simply had no place for women preachers, and he later backed her personally after she preached in public. Smith ran into something more organized. By the 1870s there was an actual faction. She overhears ministers saying ``Woman preacher; they want to be ordained'' and ``I mean to fight that thing.'' The opposition had become a position people defended on purpose.

Smith also faces exclusion Lee does not write about much. Her story is doubled. White churches push her out because she is Black, and her own denomination pushes her out because she is a woman. Being thrown out of a white church and being whispered about at an AME conference happen in the same few pages.

Together they show that the AME did not fix the problem it was founded over. Allen walked out of St. George's because white Methodists treated Black members as lesser, then built a church that did the same thing to women. The AME let women preach locally in 1884 but withheld ordination until 1948. Exclusion was essentially cloned by the solution that was supposed to stop it.

\end{document}
